%%%%%%%%%%%%%%%%%%%%%%%%%%%%%%%%%%%%%%%
% Cover Letter - IBM Client Innovation Center Germany, Data Scientist
%%%%%%%%%%%%%%%%%%%%%%%%%%%%%%%%%%%%%%%

\documentclass[]{cover}
\usepackage{fancyhdr}

\pagestyle{fancy}
\fancyhf{}

\rfoot{Page \thepage \hspace{0pt}}
\thispagestyle{empty}
\renewcommand{\headrulewidth}{0pt}
\begin{document}

%%%%%%%%%%%%%%%%%%%%%%%%%%%%%%%%%%%%%%
%     TITLE NAME
%%%%%%%%%%%%%%%%%%%%%%%%%%%%%%%%%%%%%%
\namesection{}{\Huge{Ana Lorena Ortiz Loyola}}{  \href{mailto:a01750283@tec.mx}{a01750283@tec.mx} | +49 176 62145797 |  \urlstyle{same}\href{https://linkedin.com/in/lorenasolana}{LinkedIn}
}

%%%%%%%%%%%%%%%%%%%%%%%%%%%%%%%%%%%%%%
%     MAIN COVER LETTER CONTENT
%%%%%%%%%%%%%%%%%%%%%%%%%%%%%%%%%%%%%%

\currentdate{\today}
\lettercontent{Sehr geehrte Damen und Herren,}

\lettercontent{mit großem Interesse bewerbe ich mich auf die Stelle als Data Scientist (m/w/d) beim IBM Client Innovation Center Germany GmbH. Mit praktischer Erfahrung im Aufbau und der Validierung prädiktiver Modelle in den Bereichen Gesundheitsdaten, Luftqualitätsdaten und kryptografische Algorithmen sowie einem soliden Fundament in Python, SQL und maschinellem Lernen bringe ich die analytische Tiefe mit, die für die Entwicklung von Machine-Learning-Pipelines in Ihrem Team gefragt ist.}

\lettercontent{Besonders reizt mich, dass das Client Innovation Center Projekte über Branchen wie Banking, Public Sector und Healthcare hinweg in agilen, kundennahen Teams umsetzt. Genau in diesem Zusammenspiel aus technischer Tiefe und verständlicher Kommunikation an unterschiedliche Stakeholder sehe ich meine Stärke, wie meine Erfahrung bei der Deutschen Bundesbank im Austausch mit meldenden Banken zeigt.}

\lettercontent{Konkret bringe ich mit:}

{\raggedright\fontspec[Path = OpenFonts/fonts/raleway/]{Raleway-Medium}\fontsize{11pt}{13pt}\selectfont
\begin{itemize}
    \item \textbf{Prädiktive Modellierung}: In einem Projekt zur Analyse von über 25 Jahren nationaler Gesundheitsdaten entwickelte und validierte ich Random-Forest- und neuronale Netzmodelle zur Vorhersage kritischer Outcomes sowie K-Means-Clustering zur Segmentierung von Hochrisikogruppen.
    \item \textbf{Klassische ML-Algorithmen}: Mit SVM-Modellen (linear und RBF) klassifizierte ich kritische Ozonwerte aus stadtweiten Luftqualitätsdaten mit hoher Genauigkeit.
    \item \textbf{Mathematisches Fundament}: In einem Forschungsprojekt zu gitterbasierter Kryptografie (LWE, NTRU) analysierte ich Sicherheitsparameter und Laufzeitverhalten; die Ergebnisse flossen in eine wissenschaftliche Publikation ein.
    \item \textbf{Moderne Entwicklungswerkzeuge}: Ich nutze im Alltag KI-gestützte Programmierwerkzeuge wie Claude Code zur Beschleunigung von Modell- und Pipeline-Entwicklung.
\end{itemize}\par}
\vspace{6pt}

\lettercontent{Ich arbeite strukturiert und analytisch, mit Fokus auf Sorgfalt und Nachvollziehbarkeit, kommuniziere Ergebnisse aber ebenso gerne verständlich an technische und nicht-technische Stakeholder. Diese Fähigkeit habe ich bei der Deutschen Bundesbank im Austausch mit meldenden Banken vertieft.}

\lettercontent{Ich freue mich auf die Gelegenheit, meine Erfahrung in Ihr Team einzubringen, und stehe für ein persönliches Gespräch gerne zur Verfügung.}

\begin{flushright}
\closing{Mit freundlichen Grüßen,}

\signature{Ana Lorena Ortiz Loyola}
\end{flushright}
\end{document}
